\section{Current, torque, and operating regions}

Current-vector control under simultaneous inverter voltage and current limits
is a classical route to extending the PSM operating range
\cite{morimoto1990,kim1997}; the derivation below makes its geometry explicit.

The stator-current constraint is
\begin{equation}
 \id^2+\iq^2\leq\Imax^2,\label{eq:current-circle}
\end{equation}
a disk because its quadratic matrix is the identity.  Unlike the voltage
ellipse, it is centered at the origin and isotropic.

For $p$ pole pairs, electromagnetic torque is
\begin{equation}
 \Te=\kt\left[\psif\iq+(\Ld-\Lq)\id\iq\right],
 \qquad \kt=\frac32p.\label{eq:psm-torque}
\end{equation}
Let $\dL=\Ld-\Lq$.  A constant nonzero torque satisfies
\begin{equation}
 \iq=\frac{\Te/\kt}{\psif+\dL\id}.\label{eq:torque-hyperbola}
\end{equation}
When $\dL\neq0$, translation $\tilde i_d=i_d+\psif/\dL$ reduces it to
$\tilde i_d i_q=\text{constant}$: a rectangular hyperbola.  When $\dL=0$,
the constant-torque curves are horizontal lines, not hyperbolas.  At zero
torque, the equation degenerates to the union of $i_q=0$ and
$i_d=-\psif/\dL$.

\begin{figure}[tb]
\centering
\begin{tikzpicture}
 \begin{axis}[axis main,width=.78\linewidth,height=7cm,
 xlabel={$i_d$ (normalized)},ylabel={$i_q$ (normalized)},
 xmin=-2.6,xmax=1.6,ymin=-1.8,ymax=1.8,axis equal image]
  \addplot[current limit,domain=0:360,samples=181]
    ({1.45*cos(x)},{1.45*sin(x)});
  \addplot[voltage limit,domain=0:360,samples=181]
    ({-1.0+1.8*cos(x)*cos(12)-.9*sin(x)*sin(12)},
     {.15+1.8*cos(x)*sin(12)+.9*sin(x)*cos(12)});
  \addplot[only marks,mark=*,black] coordinates {(-1.2,1.05)};
  \legend{current limit,voltage limit,feasible operating point}
 \end{axis}
\end{tikzpicture}
\caption{The feasible currents lie in the intersection of the current disk and
voltage ellipse.}
\label{fig:current-voltage}
\end{figure}

\begin{figure}[tb]
\centering
\begin{tikzpicture}
 \begin{axis}[axis main,width=.8\linewidth,height=6.6cm,
 xlabel={$i_d$},ylabel={$i_q$},xmin=-3.8,xmax=1.2,ymin=-2.4,ymax=2.4,
 restrict y to domain=-2.5:2.5]
  \foreach \t in {-1.2,-.6,.6,1.2}{
   \addplot[torque contour,domain=-3.7:1.1,samples=160]
     {\t/(1-.45*x)};
  }
  \addplot[densely dashed,gray] coordinates {(2.222,-2.4) (2.222,2.4)};
  \addplot[densely dashed,gray] coordinates {(-3.8,0) (1.2,0)};
 \end{axis}
\end{tikzpicture}
\caption{Constant nonzero torque curves are shifted hyperbolas when saliency is
nonzero.  The zero-torque level degenerates into its asymptote lines.}
\label{fig:torque-contours}
\end{figure}


\subsection{Tangency interpretation}

\gls{mtpa} minimizes $\norm{\current}^2$ at fixed torque.  The multiplier
condition is best discovered from the picture rather than memorized as a
recipe.  A current circle is a level set of
$f(\current)=\norm{\current}^2$; a torque contour is a level set of
$g(\current)=\Te(\current)$.  At the smallest circle that reaches the requested
torque, the two curves touch.  Their tangents coincide, so their normals are
parallel.  Because a gradient is normal to a regular level set, there must be
a scalar $\mu$ such that
\begin{equation}
 2\current=\mu\nabla_{\current}\Te.
\end{equation}
This geometric necessity is the Lagrange multiplier condition.  It exposes
the optimum without parameterizing the hyperbola, but it does not by itself
select the correct branch or guarantee global optimality; feasibility and
comparison of stationary candidates remain necessary.
Eliminating $\mu$ gives the familiar locus
\begin{equation}
 \psif\id+\dL(\id^2-\iq^2)=0.\label{eq:mtpa-locus}
\end{equation}
Geometrically, the smallest current circle reaching a requested torque contour
is tangent to it.

As speed rises, the voltage ellipse contracts and moves toward the flux
cancellation point.  \gls{fw} begins when the unconstrained MTPA point no
longer satisfies the voltage inequality.  With both current and voltage active,
the feasible operating point lies at their intersection and the KKT condition
expresses the torque normal as a nonnegative combination of both constraint
normals \cite{boyd2004}.

The nonnegative coefficients in this combination are not arbitrary algebraic
decorations.  Each active inequality can oppose motion only toward its
infeasible side.  KKT therefore replaces one parallel-normal condition by a
cone of admissible active-boundary normals.  An inactive boundary has zero
multiplier and contributes no normal.

In the ideal $R_s=0$ model, \gls{mtpv} maximizes torque on
\begin{equation}
 \omegae^2\left[\Lq^2\iq^2+(\Ld\id+\psif)^2\right]=\Umax^2.
\end{equation}
It is again a tangency, now between a torque contour and the voltage ellipse.
Using the gradients of torque and ideal voltage and eliminating the Lagrange
multiplier yields the explicit locus
\begin{equation}
 \boxed{\dL\Lq^2\iq^2
 =\Ld(\psif+\dL\id)(\psif+\Ld\id).}\label{eq:ideal-mtpv}
\end{equation}
This relation is valid on regular branches for $R_s=0$; feasibility and the
sign of torque must still be checked.  For $\dL=0$ it reduces to the appropriate
SPMSM limiting condition rather than a hyperbolic saliency branch.
The term is meaningful only after current or thermal admissibility and the
chosen machine model have been stated.

\begin{figure}[tb]
\centering
\begin{tikzpicture}
 \begin{axis}[axis main,width=.82\linewidth,height=7cm,
 xlabel={$i_d$},ylabel={$i_q$},xmin=-2.5,xmax=.5,ymin=0,ymax=1.9,
 legend pos=north west]
  \addplot[current limit,domain=0:180,samples=120]
    ({1.55*cos(x)},{1.55*sin(x)});
  \addplot[voltage limit,domain=0:360,samples=181]
    ({-1.18+1.35*cos(x)},{.72*sin(x)});
  \addplot[optimization path,domain=-1.65:-.08,samples=80]
    ({x},{sqrt(max(0,x*x+x/(-.38)))});
  \addplot[only marks,mark=*,mark size=2pt] coordinates {(-.34,1.05) (-1.22,1.0) (-1.7,.58)};
  \node[annotation] at (axis cs:-.15,.75){MTPA};
  \node[annotation] at (axis cs:-1.05,1.6){FW};
  \node[annotation] at (axis cs:-2.05,.32){MTPV side};
 \end{axis}
\end{tikzpicture}
\caption{Conceptual active-set geometry: MTPA is current-limited, field
weakening uses both limits, and the high-speed branch follows the voltage
boundary toward an MTPV tangency.}
\label{fig:operating-regions}
\end{figure}


\paragraph{Synthesis.}
MTPA, field weakening, and MTPV are not three unrelated formulas.  They are
successive contact configurations of torque level sets with a current circle,
then a voltage ellipse, and sometimes both.  The algebra changes because the
active normals change; the geometric principle does not.
